\documentclass[../../main.tex]{subfiles}
\graphicspath{{../../image/}{../../figure/}}
\begin{document}

\section{Young 测度}\label{section:ch17-young测度}

构造格式的相容性之所以难以确立，关键在于生成近似解弱收敛序列的那些格式，一般不允许在非线性函数下取弱极限。设 $\Omega$ 是 $\mathbb{R}^m$ 的开子集，$\{U_k\}$ 是 $L^\infty(\Omega;\mathbb{R}^n)$ 中在 $L^\infty$ 弱$^*$ 意义下收敛于某极限 $\bar U$ 的序列。若 $g$ 是 $\mathbb{R}^n$ 上任意连续实值函数，则序列 $\{g(U_k)\}$ 含有在 $L^\infty$ 弱$^*$ 意义下收敛的子列，譬如收敛于 $\bar g$，但一般 $\bar g\neq g(\bar U)$。事实表明，对一切连续 $g$，这类序列的极限行为被一族定义在 $\mathbb{R}^n$ 上的概率测度 $\{\nu_X:X\in\Omega\}$ 所编码，该族由下述步骤构造。

记 $M(\mathbb{R}^n)$ 为 $\mathbb{R}^n$ 上有界 Radon 测度空间，它等距同构于有界连续函数空间 $C(\mathbb{R}^n)$ 的对偶。对 $k=1,2,\cdots$ 以及任意 $X\in\Omega$，我们联系一个集中于点 $U_k(X)$ 的 Dirac 质量 $\delta_{U_k(X)}\in M(\mathbb{R}^n)$，并把族 $\{\delta_{U_k(X)}:X\in\Omega\}$ 实现为空间 $L^\infty_w(\Omega;M(\mathbb{R}^n))$ 中的一个元素 $\nu_k$，该空间等距同构于 $L^1(\Omega;C(\mathbb{R}^n))$ 的对偶。由标准的弱紧性与可分性定理，存在 $\{\nu_k\}$ 的子列 $\{\nu_j\}$，它弱$^*$ 收敛于某个 $\nu\in L^\infty_w(\Omega;M(\mathbb{R}^n))$。于是 $\nu=\{\nu_X:X\in\Omega\}$，且当 $j\to\infty$ 时
\begin{align}
\int_\Omega\psi(X,U_j(X))\,dX=\int_\Omega\langle\delta_{U_j(X)},\psi(X,\cdot)\rangle\,dX\longrightarrow\int_\Omega\langle\nu_X,\psi(X,\cdot)\rangle\,dX,
\label{align:ch17-1711}
\end{align}
对任意 $\psi\in C(\Omega\times\mathbb{R}^n)$ 成立。$\delta_{U_j(X)}$ 的支撑一致有界，因而 $\nu_X$ 必具有紧支撑。此外，由于 $\delta_{U_j(X)}$ 是概率测度，$\nu_X$ 也是概率测度。特别地，在 $\psi(X,U)=\phi(X)g(U)$（其中 $\phi\in C(\Omega)$，$g\in C(\mathbb{R}^n)$）的情形下应用\eqref{align:ch17-1711}，我们得到下述

\begin{theorem}[Young 测度的存在性]\label{theorem:ch17-young测度存在性}
设 $\Omega$ 是 $\mathbb{R}^m$ 的开子集。则 $L^\infty(\Omega;\mathbb{R}^n)$ 中任意有界序列 $\{U_k\}$ 都含有一个子列 $\{U_j\}$ 以及一个具有紧支撑的概率测度可测族 $\{\nu_X:X\in\Omega\}$，使得对任意 $g\in C(\mathbb{R}^n)$，
\begin{align}
g(U_j)\rightharpoonup\bar g,\qquad j\to\infty,
\label{align:ch17-1712}
\end{align}
在 $L^\infty$ 弱$^*$ 意义下成立，其中
\begin{align}
\bar g(X)=\langle\nu_X,g\rangle=\int_{\mathbb{R}^n}g(U)\,d\nu_X(U).
\label{align:ch17-1713}
\end{align}
\end{theorem}

族 $\{\nu_X:X\in\Omega\}$ 构成与子列 $\{U_j\}$ 相联系的 Young 测度族。为获得直观认识，考虑 $\Omega$ 中以某点 $X\in\Omega$ 为心、半径为 $r$、测度为 $|\mathscr{B}_r|$ 的球 $\mathscr{B}_r(X)$。由 $\nu_X$ 的构造容易看出
\begin{align}
\nu_X=\lim_{r\downarrow0}\lim_{j\uparrow\infty}\frac{1}{|\mathscr{B}_r|}\int_{\mathscr{B}_r(X)}\delta_{U_j(Y)}\,dY,\qquad\text{a.e. on }\Omega,
\label{align:ch17-1714}
\end{align}
其中极限按弱$^*$ 意义理解。注意\eqref{align:ch17-1714}右端的平均积分可解释为：当 $Y$ 从 $\mathscr{B}_r(X)$ 中均匀随机抽取时，$U_j(Y)$ 取值的概率分布。于是由\eqref{align:ch17-1714}，$\nu_X$ 表示 $X$ 附近 $U_j$ 取值的极限概率分布。

由\eqref{align:ch17-1712}、\eqref{align:ch17-1713}，子列 $\{U_j\}$ 在 $L^\infty$ 弱$^*$ 意义下收敛于 Young 测度的均值 $\bar U=\langle\nu_X,U\rangle$。对一切 $g\in C(\mathbb{R}^n)$，$\{g(U_j)\}$ 的极限 $\bar g$ 满足 $\bar g=g(\bar U)$，当且仅当 $\nu_X$ 退化为集中于 $\bar U(X)$ 的 Dirac 质量 $\delta_{\bar U(X)}$。在这种情形下，$\{|U_j|\}$ 将收敛于 $|\bar U|$，从而 $\{U_j\}$ 在 $L^p_{\mathrm{loc}}(\Omega)$ 中对任意 $1\leqslant p<\infty$ 强收敛于 $\bar U$，且 $\{U_j\}$ 的某个子列在 $\Omega$ 上 a.e. 收敛于 $\bar U$。因此，要确立 $\{U_j\}$ 的强收敛，就需要验证 Young 测度的支撑被限制在一点上。

某些应用要求\cref{theorem:ch17-young测度存在性}更一般的版本。即使序列 $\{U_k\}$ 只在某个 $L^p(\Omega;\mathbb{R}^n)$（$1<p<\infty$）中有界，Young 测度 $\nu_X$ 仍可定义。若 $\Omega$ 有界，则 $\nu_X$ 仍是概率测度，且\eqref{align:ch17-1712}、\eqref{align:ch17-1713}对一切满足增长条件 $|g(U)|\leqslant c(1+|U|^q)$（其中 $0<q<p$）的连续函数 $g$ 成立。此时\eqref{align:ch17-1712}中的收敛在 $L^r(\Omega)$ 中为弱收敛，其中 $1<r<p/q$。相反，当 $\Omega$ 无界时，$\nu_X$ 的质量可能小于一，因为在构造它们的过程中，当取 $j\to\infty$ 极限时，部分质量可能在无穷远处泄漏。

\end{document}
